\documentclass[11pt]{article}

\usepackage[utf8]{inputenc}
\usepackage[T1]{fontenc}
\usepackage[margin=1in]{geometry}
\usepackage{booktabs}
\usepackage{longtable}
\usepackage{xcolor}
\usepackage{hyperref}
\usepackage{enumitem}
\usepackage{array}
\usepackage{parskip}

\hypersetup{colorlinks=true, linkcolor=blue, urlcolor=blue, citecolor=blue}

\title{Retrieval and Classification of the Scientific Literature on the\\
Capacitated Arc Routing Problem: A Reproducible Pipeline}
\author{CARP review project}
\date{\today}

\begin{document}
\maketitle

\begin{abstract}
This document describes the methodology used to systematically collect, classify, and
analyse the scientific literature on the Capacitated Arc Routing Problem (CARP), its
structural variants, and its solution methodology. The pipeline has three stages:
(1) automated metadata retrieval from the OpenAlex bibliographic aggregator,
(2) rule-based classification of each article's CARP variant and solution approach
through a structured keyword taxonomy, and (3) semantic classification with a large
language model (LLM) that assigns a single primary label per article. The corpus
covers 905 unique works (1981--2026); a focused subset of 206 works published from
2021 onward is used for the main analysis.
\end{abstract}

% ------------------------------------------------------------------ %
\section{Data collection}
% ------------------------------------------------------------------ %

The pipeline deliberately avoids scraping publisher websites (Elsevier ScienceDirect,
IEEE~Xplore, Springer Link, Wiley), which is technically brittle---anti-bot
protections, CAPTCHAs, dynamic rendering---and prohibited by most terms of service.
Instead it queries \textbf{OpenAlex} (\url{https://openalex.org}), a free, openly
licensed catalogue of more than 250 million scholarly works aggregating Crossref,
PubMed, DOAJ, Unpaywall, and all major publishers through a public REST API. Only
metadata are stored; no full texts or PDFs are downloaded. Supplying a contact
e-mail places the client in the faster ``polite pool'' of the API.

OpenAlex distributes abstracts as an \emph{inverted index}---a mapping from each
distinct word to its positions in the original text---for licensing reasons. The
readable abstract is reconstructed by sorting all (position, word) pairs by position
and concatenating the words. The procedure is lossless for vocabulary and word order.
When OpenAlex holds no abstract for a work, the record is kept in the metadata corpus
and in all year-based counts, but is excluded from the text-based classification
stages, which require readable text.

\subsection{Search strategy}

The search issues 22 exact-phrase queries against the combined title-and-abstract
field, organised in three groups:

\begin{itemize}[nosep]
  \item \textbf{Core CARP (3):} \emph{capacitated arc routing problem};
        \emph{capacitated arc routing}; \emph{arc routing problem}.
  \item \textbf{Structural variants (11):} \emph{periodic capacitated arc routing};
        \emph{stochastic arc routing}; \emph{mixed capacitated arc routing};
        \emph{multi-depot arc routing}; \emph{multi-objective arc routing};
        \emph{split delivery arc routing}; \emph{time-dependent arc routing};
        \emph{min-max capacitated arc routing}; \emph{large-scale capacitated arc
        routing}; \emph{windy arc routing}; \emph{open capacitated arc routing}.
  \item \textbf{Methodology (8):} \emph{arc routing} followed by
        \emph{metaheuristic}; \emph{heuristic}; \emph{genetic algorithm};
        \emph{memetic algorithm}; \emph{ant colony}; \emph{tabu search};
        \emph{exact algorithm}; \emph{branch and cut}.
\end{itemize}

Results from all queries are pooled and de-duplicated by the persistent OpenAlex work
identifier, so each article appears at most once. Pagination follows the API's
cursor scheme (pages of up to 200 works, chained via a \emph{next-cursor} token)
until each result set is exhausted, with a courtesy delay between requests.

For every unique work the following fields are stored: DOI (normalised), title,
reconstructed abstract, publication year, venue, publisher, open-access flag and
route (gold/green/hybrid/bronze/closed), work type, cited-by count, and the OpenAlex
identifier.

\subsection{Collection results}

The unfiltered query returned \textbf{905 unique works} dated 1969--2026. Records
predating the formal introduction of the CARP by Golden and Wong (1981) correspond to
a few mis-dated OpenAlex entries and to predecessor problems (Chinese and Rural
Postman) that share the arc-routing vocabulary. Approximately 79\% of records carry a
DOI, 61\% include a usable abstract, and 35\% are open access; missing abstracts
reflect coverage gaps in the source data rather than retrieval failures.

Because the pre-2021 portion contains many abstract-less and tangential records, the
main analysis focuses on the \textbf{2021--2026 subset}, obtained by filtering the
already-collected data without re-querying or re-classifying. It contains
\textbf{206 works} (44, 55, 30, 26, 31, and 20 per year from 2021 to 2026), of which
135 carry a usable abstract and 113 are open access.

% ------------------------------------------------------------------ %
\section{Rule-based classification}
% ------------------------------------------------------------------ %

The first classification stage is a deterministic keyword tagger: the title and
abstract of each article are concatenated and matched against four taxonomies of
word-boundary regular expressions (case-insensitive, except for all-capital
acronyms). Tables~\ref{tab:variants}--\ref{tab:metas} enumerate every category and
its trigger keywords.

\subsection{CARP variant taxonomy}

Nineteen structural categories are recognised (Table~\ref{tab:variants}). A category
is assigned whenever any of its keywords appears; multiple categories may co-occur
(e.g.\ a multi-depot CARP with time windows). If no keyword is detected, the article
is labelled \emph{Classical CARP / unspecified}.

{\small
\begin{longtable}{@{}p{4.3cm}p{9.3cm}@{}}
\caption{CARP variant taxonomy and trigger keywords.}\label{tab:variants}\\
\toprule
\textbf{Category} & \textbf{Trigger keywords / acronyms} \\
\midrule
\endfirsthead
\toprule
\textbf{Category} & \textbf{Trigger keywords / acronyms} \\
\midrule
\endhead
Periodic (PCARP) &
  ``periodic (capacitated) arc routing''; acronym PCARP. \\
Stochastic &
  ``stochastic (capacitated) arc routing''; ``stochastic demand''; acronym SCARP. \\
Uncertain / Fuzzy / Robust &
  ``fuzzy''; ``robust''; ``uncertainty''; ``chance-constrain''. \\
Mixed (MCARP) &
  ``mixed (capacitated) arc routing''; acronym MCARP. \\
Multi-depot &
  ``multi-depot'' (with or without hyphen); ``multiple depots''. \\
Multi-objective &
  ``multi-objective''; ``bi-objective''; ``pareto''. \\
Min-max / Balanced &
  ``min-max''; any form of ``balanc'' (balance, balanced, balancing). \\
Split delivery &
  ``split-delivery''; ``split-demand''. \\
Time-dependent / Dynamic &
  ``time-dependent arc routing''; ``dynamic arc routing''; ``online arc routing''. \\
Time windows (CARPTW) &
  ``time window''; acronym CARPTW. \\
Windy (WARP) &
  ``windy rural postman''; ``windy arc routing''; ``windy postman''; ``windy''. \\
Open (OCARP) &
  ``open (capacitated) arc routing''; acronym OCARP. \\
Prize-collecting / Profitable &
  ``prize-collecting''; ``profitable''; ``team orienteering arc''. \\
Large-scale &
  ``large-scale'' (with or without hyphen). \\
Intermediate facilities (IF) &
  ``intermediate facilit''; ``multiple trips''; ``replenish''; acronym IF (not
  followed by a letter). \\
Green / Electric &
  ``green'' (as adjective); ``electric vehicle''; ``emission''; ``CO2'';
  ``low-carbon''. \\
Node / General routing (GRP) &
  ``general routing problem''; ``node, edge'' (combination); acronym GRP. \\
Rural postman (RPP) &
  ``rural postman''; acronym RPP. \\
Chinese postman (CPP) &
  ``chinese postman''; acronym CPP. \\
\bottomrule
\end{longtable}
}

\subsection{Method taxonomies}

Eight exact-method families (Table~\ref{tab:exact}), six constructive-heuristic
families (Table~\ref{tab:constructive}), and seventeen metaheuristic families
(Table~\ref{tab:metas}) are detected, in most cases via both the full name and the
widely used acronym.

{\small
\begin{longtable}{@{}p{4.3cm}p{9.3cm}@{}}
\caption{Exact-method taxonomy.}\label{tab:exact}\\
\toprule
\textbf{Category} & \textbf{Trigger keywords} \\
\midrule
\endfirsthead
\toprule
\textbf{Category} & \textbf{Trigger keywords} \\
\midrule
\endhead
Branch-and-cut & ``branch-and-cut''; ``branch-cut''. \\
Branch-and-price & ``branch-and-price''; ``branch-price''. \\
Branch-and-bound & ``branch-and-bound''; ``branch-bound''. \\
Column generation & ``column generation''. \\
Cutting plane & ``cutting-plane''. \\
Integer / linear programming &
  ``integer (linear) program''; ``linear program''; ``mathematical program'';
  acronyms MILP, MIP, ILP. \\
Dynamic programming & ``dynamic programming''. \\
Lagrangian relaxation & ``lagrang'' (Lagrangian, Lagrange relaxation, etc.). \\
\bottomrule
\end{longtable}
}

{\small
\begin{longtable}{@{}p{4.3cm}p{9.3cm}@{}}
\caption{Constructive-heuristic taxonomy.}\label{tab:constructive}\\
\toprule
\textbf{Category} & \textbf{Trigger keywords} \\
\midrule
\endfirsthead
\toprule
\textbf{Category} & \textbf{Trigger keywords} \\
\midrule
\endhead
Path-scanning & ``path-scanning''. \\
Augment-merge & ``augment-merge''; ``augment-insert''. \\
Construct-strike & ``construct-strike''. \\
Ulusoy split / route-first & ``ulusoy''; ``route-first''; ``split procedure''. \\
Savings algorithm &
  ``savings algorithm''; ``savings heuristic''; ``savings-based''. \\
Greedy / Nearest neighbour & ``greedy''; ``nearest-neighbour''. \\
\bottomrule
\end{longtable}
}

{\small
\begin{longtable}{@{}p{5.0cm}p{8.6cm}@{}}
\caption{Metaheuristic taxonomy.}\label{tab:metas}\\
\toprule
\textbf{Category} & \textbf{Trigger keywords / acronyms} \\
\midrule
\endfirsthead
\toprule
\textbf{Category} & \textbf{Trigger keywords / acronyms} \\
\midrule
\endhead
Genetic / Evolutionary algorithm &
  ``genetic algorithm''; ``evolutionary algorithm''; ``evolutionary computation'';
  ``genetic programming''; acronym GA. \\
Memetic algorithm & ``memetic''. \\
Ant colony optimisation (ACO) & ``ant colony''; ``ant system''; acronym ACO. \\
Tabu search & ``tabu search''; ``tabu''. \\
Simulated annealing & ``simulated annealing''. \\
Variable neighbourhood search (VNS) &
  ``variable neighbourhood'' (British or American spelling); acronyms VNS, VND. \\
Particle swarm optimisation (PSO) & ``particle swarm''; acronym PSO. \\
GRASP & ``greedy randomized adaptive''; acronym GRASP. \\
Iterated local search (ILS) & ``iterated local search''; acronym ILS. \\
Guided local search & ``guided local search''. \\
Scatter search / Path relinking & ``scatter search''; ``path-relinking''. \\
Estimation of distribution (EDA) & ``estimation of distribution''; acronym EDA. \\
Artificial bee colony (ABC) &
  ``artificial bee colony''; ``bee colony''; ``ABC algorithm''. \\
Differential evolution (DE) & ``differential evolution''. \\
Large neighbourhood search (LNS/ALNS) &
  ``large neighbourhood'' (British or American spelling); acronyms ALNS, LNS. \\
Hyper-heuristic & ``hyper-heuristic''. \\
Cooperative / Swarm (other) &
  ``firefly''; ``cuckoo''; ``grey wolf''; ``harmony search''; ``water flow'';
  ``immune algorithm''; ``bacterial''; acronym ABC (general usage). \\
\bottomrule
\end{longtable}
}

\subsection{Solution-approach assignment and hybrid detection}

Each article receives a single high-level \emph{solution approach} label following a
priority hierarchy:

\begin{enumerate}[nosep]
  \item \textbf{Hybrid / Matheuristic} --- an explicit hybrid signal is present
        (``hybrid'', ``matheuristic'', ``math-heuristic'', ``cooperative'',
        ``memetic''), \emph{or} at least one exact method co-occurs with at least one
        metaheuristic, \emph{or} two or more distinct metaheuristic families are
        detected simultaneously.
  \item \textbf{Metaheuristic} --- at least one metaheuristic family, not hybrid.
  \item \textbf{Constructive heuristic} --- at least one constructive family, none of
        the above.
  \item \textbf{Exact} --- at least one exact method, none of the above.
  \item \textbf{Generic heuristic} --- only generic vocabulary (``heuristic'',
        ``local search'', ``neighbourhood search'', ``metaheuristic'') with no
        specific family identified.
  \item \textbf{Unspecified} --- no method-related keyword found.
\end{enumerate}

Because the tagger matches every keyword present in the text, it may over-label
survey-style abstracts that enumerate many methods; it is therefore treated as a
fast, transparent first pass rather than the sole source for precision
classification.

% ------------------------------------------------------------------ %
\section{LLM semantic classification}
% ------------------------------------------------------------------ %

To overcome the over-labelling of the keyword tagger, every article with a usable
abstract---552 of the 905 works (61\%)---is re-classified by a large language model
(Anthropic Claude). The model receives the title and the first 800~characters of the
abstract, and, unlike the rule-based pass, assigns \emph{one} primary variant and
\emph{one} primary approach per article using the closed six-field schema of
Table~\ref{tab:schema}. Its main advantages are semantic inference (methods described
by paraphrase are recognised without exact keywords), explicit relevance filtering of
records that are not about arc routing, and a one-sentence summary of how each
problem was solved. Records labelled \emph{Not relevant} are excluded from all
thematic analyses.

{\small
\begin{longtable}{@{}p{3.4cm}p{10.2cm}@{}}
\caption{LLM classification schema.}\label{tab:schema}\\
\toprule
\textbf{Field} & \textbf{Description and permitted values} \\
\midrule
\endfirsthead
\toprule
\textbf{Field} & \textbf{Description and permitted values} \\
\midrule
\endhead
Relevance &
  \emph{Core CARP}; \emph{CARP variant}; \emph{Related arc routing} (RPP, CPP,
  Windy Postman, etc.); \emph{Not relevant}. \\
Variant &
  A single label from a closed vocabulary of CARP variants (Classical CARP,
  Periodic, Stochastic, Multi-depot, Time windows, Green/Electric, \ldots) that best
  describes the paper's primary contribution. \\
Approach &
  A single label: \emph{Exact}; \emph{Constructive heuristic}; \emph{Metaheuristic};
  \emph{Hybrid / Matheuristic}; \emph{Machine learning};
  \emph{Survey / Theoretical / Modelling}. \\
Metaheuristics &
  List of the specific metaheuristic families named in the text; empty when the
  approach is not metaheuristic. \\
Is hybrid &
  Boolean: true when two or more distinct algorithmic families are combined. \\
How solved &
  One English sentence summarising, at an abstract level, how the problem was
  solved. \\
\bottomrule
\end{longtable}
}

% ------------------------------------------------------------------ %
\section{Analysis, reproducibility, and limitations}
% ------------------------------------------------------------------ %

From the LLM-classified corpus the following outputs are produced, for the full
corpus and independently for the 2021-onward subset: publications per year; variant
and solution-approach distributions; the 15 most frequent metaheuristic families;
open-access share; the yearly composition of approaches from 2005 onward; and a
variant~$\times$~approach cross-tabulation exported as CSV.

The pipeline is reproducible from any internet-connected machine: it queries only a
public API, involves no manual curation, and every transformation is deterministic
except the LLM call (non-deterministic but statistically stable). All intermediate
files are version-controlled, so figures and tables can be regenerated without
re-querying OpenAlex.

Three limitations apply. \emph{Coverage:} OpenAlex does not index every venue; the
39\% of records without an abstract mostly reflects this gap. \emph{Classifier:} the
rule-based tagger operates on surface text only, missing non-standard terminology
and over-labelling surveys; the LLM mitigates this but is limited to records with an
abstract. \emph{Temporal scope:} the API was queried at a fixed point in time;
re-running the collection and de-duplicating against existing identifiers would
extend the corpus without discarding prior classifications.

\end{document}
